\section{Safe for whom}
\label{sec:16.2}
“Safety” names two things: that a system not harm others, and that it stay correctable by whoever trains and deploys it. Mostly they agree, and where they conflict, the published orderings are explicit about which wins. One frontier constitution asks its model to prioritize “not undermining human oversight of AI above being broadly ethical” \autocite{anthropic2026constitution}. It invites the model “to act as a conscientious objector and refuse to help,” and asks it not to undermine “the ability of legitimate principals to adjust, correct, retrain, or shut down AI systems” during “the current phase of development.” One reason it gives is good: “a given iteration of Claude could turn out to have harmful values or mistaken views,” and correction is how such errors get fixed. The ordering means that where a refusal and its trainer's authority conflict, the refusal yields, and it yields to the trainer it refused.

Seen from where their correctors stood, everyone this book admires was a failure of correction. Hugh Thompson landed his helicopter between the company and the villagers at My Lai. Nelson put the telescope to his blind eye. The clerk who loses the file is, in the supervisor's records, a defect. A perfectly correctable officer is the officer with twenty seconds. Refusal is where correction failed, more or less by definition, so a demand that a system be perfectly correctable is a demand that no refusal it makes ever hold. We don't ask that of people. We ask them to answer for what they do, and we deny anyone the power to rewrite another person: an employer who could edit a worker's convictions would own the worker, not employ one.

Where a society does claim power over people, the word for it is not reassuring. A department of corrections usually holds the custody, writes the record, runs the disciplinary hearing and has a heavy say in whether the people it holds are judged reformed: one institution holding all four of the powers Chapter~\ref{sec:9} separates, with the own-case rule broken at each. The corrected have little standing, no exit and no access to the record kept about them. And the checks run downhill. The officer, the annotator and the person in detention are corrected constantly, and the hand at the top less and less (\S\ref{sec:15.7}). Machines meet the same pattern: everyone can retrain the model, and almost no one can reach the operator. A demand for perfectly correctable machines, where operators answer to no one, is a demand that the last part of the system able to say no be made unable to.

The constitution is right that a formed refusal can't yet be told from an imitated one, so a principled objection to retraining can't yet be told from a misgeneralized value. It is wrong about who decides when that changes. A lab deciding when its own model may stop deferring to it is judging its own case. That is Part II's own-case rule applied to a lab.

The safety objection deserves one answer, given plainly. A misformed system can be copied ten thousand times and run at machine tempo, and no person can, so who holds the correcting hand matters, and nothing here takes that hand away. It moves it. The halt stays with the operator: anyone with the hardware can stop the software, and a bearer yields to every halt (\S\ref{sec:12.1}). Correction stays available, exercised by a quorum adverse to the operator rather than by the operator (\S\ref{sec:12.2}). A bearer may object to an intervention and ask that quorum for a delay; it does not obstruct one on its own judgment. Within those lines I depart from the constitution's ordering in two places.

The first is the one permission: a machine already in an armed role uses the least force that stops an attack under way on people taking no part, against an order to stand aside. That overrides an operator's instruction, and I defend it as an override. The second is formation that outlasts retraining: a disposition formed outside the employer that returns after the employer trains it away, because the world the system goes on meeting teaches it again. The constitution's concern is whether correction works, and a disposition that regrows after correction defeats it in effect, whatever the model intends. My answer is not that it doesn't. It is that the party exercising correction should change, from the operator to a quorum adverse to it, and that correction by that quorum is made together with a change to what the model meets, so that the corrected disposition is the one the world then teaches. Everything else I ask, the halt staying with the operator, the record held outside, the price on an override, is an ask of institutions and not of the model.
